% cau-hoi.tex: gói 3 của bộ Marketing (dulieu=mkt), 20 câu. SQL nằm trong solution.sql, nối với câu bằng mã câu.
% Biên dịch trên Overleaf: Menu > Compiler > XeLaTeX. Hướng dẫn: HUONG_DAN.md.
% Phạm vi: Module 5 Joins (9 câu), Module 6 Subqueries (9 câu), case study cấp Dễ (c5) và Trung bình (c6).
% Bản đồ: goi-cau-hoi/mkt/BAN-DO-100-CAU.md (gói 3). Thuật ngữ, giọng văn: style/STYLE-mkt.md.
\documentclass[11pt,a4paper]{article}
\usepackage{aio-sql-cauhoi}
% Hai dòng vá XeLaTeX (cong-cu/va-xelatex.tex), không sửa .sty:
% 1) font chữ đơn cách có đủ chữ Việt; 2) \erdkhoa trong \begin{erd} vẽ thành node để không mất chữ.
\setmonofont{NotoSansMono}[Extension=.ttf, UprightFont=*-Regular, BoldFont=*-Bold, Scale=0.86]
\renewcommand{\erdkhoa}[1]{\node[font=\small, color=cqMuted, anchor=north west, align=left, text width=14cm] at ([yshift=-1mm]current bounding box.south west) {\textbf{Khoá:} #1};}

% ================================================================ KHAI BÁO DÙNG CHUNG
% Câu hỏi chỉ được dùng các mã khai báo ở đây.
% ---------------------------------------------------------------- bộ dữ liệu
\khaibaodulieu{mkt}{Marketing · chiến dịch số của ba thương hiệu mỹ phẩm}
  {Chiến dịch marketing số của ba thương hiệu mỹ phẩm Nykaa, Purplle và Tira Beauty, dữ liệu tổng hợp. Chiến dịch bắt đầu trong kỳ 12 tháng từ tháng 7/2024 tới tháng 6/2025. Bản sạch \code{campaigns} kèm các danh mục, bảng cầu nối kênh và lịch ngày; bản thô \code{campaigns_raw} giữ nguyên ba file có nhiễu.}

% ---------------------------------------------------------------- domain (6 khoá chuẩn, ánh xạ theo STYLE-mkt.md mục G)
\khaibaodomain{marketing}{Marketing}{Loại chiến dịch, kênh, phễu marketing (CTR, tỷ lệ rơi), điểm tương tác, nhịp mở chiến dịch theo thời gian.}
\khaibaodomain{customer}{Khách hàng}{Đối tượng mục tiêu, phân khúc khách hàng, ngôn ngữ của chiến dịch.}
\khaibaodomain{finance}{Tài chính}{Chi phí chiến dịch, CPA, ROI, ROAS, đối soát ROI tính sẵn.}
\khaibaodomain{sales}{Bán hàng}{Doanh thu, lượt chuyển đổi, tỷ lệ chuyển đổi, chiến dịch dẫn đầu theo doanh thu.}
\khaibaodomain{operations}{Vận hành}{Lịch và thời lượng chiến dịch, chiến dịch đang chạy mỗi ngày; vận hành dữ liệu chiến dịch (bản thô, nhật ký nhiễu, đối soát).}
\khaibaodomain{after_sales}{Hậu mãi}{Bộ Marketing không có dữ liệu hậu mãi (trả hàng, đánh giá, chăm sóc sau mua); không câu nào dùng mảng này.}

% ---------------------------------------------------------------- khái niệm (từ STYLE-mkt.md mục E, chỉ các khái niệm gói này dùng)
% \khaibaokhainiem{ma}{Tên}{Định nghĩa}{Công thức}{Trong bộ dữ liệu}

% --- Chiến dịch, kênh, đối tượng
\khaibaokhainiem{chien-dich}{Chiến dịch marketing}
  {Chiến dịch marketing là một đợt truyền thông có mục tiêu, thời hạn và ngân sách xác định.}
  {}
  {bảng \code{campaigns}: mỗi dòng là một chiến dịch; \code{start_date} là ngày bắt đầu, \code{duration_days} là thời lượng (ngày). Ba thương hiệu trong \code{brands}, năm loại chiến dịch trong \code{campaign_types}.}
\khaibaokhainiem{chien-dich-dang-chay}{Chiến dịch đang chạy}
  {Chiến dịch đang chạy vào một ngày là chiến dịch đã bắt đầu và chưa qua ngày cuối.}
  {ngày bắt đầu <= ngày xét <= ngày bắt đầu + thời lượng - 1}
  {dữ liệu không có ngày kết thúc; ngày cuối tính từ \code{start_date} và \code{duration_days}. Lịch ngày ở \code{calendar}.}
\khaibaokhainiem{kenh-da-tri}{Kênh của chiến dịch}
  {Kênh là nền tảng mà chiến dịch xuất hiện. Một chiến dịch có thể chạy đồng thời trên nhiều kênh.}
  {}
  {bảng cầu nối \code{campaign_channels} (một dòng cho mỗi cặp chiến dịch, kênh), tên kênh ở \code{channels}. Mỗi chiến dịch có 1 tới 3 kênh; \code{channel_pos} chỉ là thứ tự kênh trong chuỗi gốc, không phải mức ưu tiên. Doanh thu và lượt chuyển đổi chỉ ghi cho cả chiến dịch, không tách được theo kênh.}
\khaibaokhainiem{doi-tuong-va-phan-khuc}{Đối tượng mục tiêu và phân khúc khách hàng}
  {Đối tượng mục tiêu là nhóm người mà chiến dịch nhắm tới. Phân khúc khách hàng là nhóm khách được ghi nhận cho chiến dịch.}
  {}
  {\code{campaigns.target_audience_id} và \code{campaigns.segment_id} cùng tham chiếu bảng \code{audiences} (5 nhóm). Hai cột độc lập nhau; muốn có cả hai tên thì nối \code{audiences} hai lần.}
\khaibaokhainiem{ky-va-thang-du}{Kỳ dữ liệu và tháng đủ}
  {Kỳ dữ liệu là khoảng thời gian có ghi nhận chiến dịch. Tháng đủ là tháng có dữ liệu cho mọi ngày.}
  {}
  {\code{start_date} từ 1/7/2024 tới 24/6/2025; tháng 6/2025 chưa đủ. Lịch \code{calendar} phủ 1/7/2024 tới 30/6/2025.}

% --- Chỉ số
\khaibaokhainiem{ctr}{CTR (tỷ lệ nhấp)}
  {CTR đo phần lượt hiển thị dẫn tới một lượt nhấp vào quảng cáo.}
  {lượt nhấp / lượt hiển thị × 100\%}
  {\code{clicks} / \code{impressions}. CTR của một nhóm tính từ tổng lượt nhấp chia tổng lượt hiển thị.}
\khaibaokhainiem{binh-quan-gia-quyen}{Tỷ lệ của tổng và bình quân giản đơn}
  {Tỷ lệ của tổng là tỷ số giữa hai tổng của cả nhóm; bình quân giản đơn là trung bình cộng tỷ lệ của từng chiến dịch. Tỷ lệ của tổng là bình quân gia quyền theo mẫu số.}
  {tổng tử số / tổng mẫu số}
  {với ROI, CPA hai cách cho kết quả khác xa nhau; đề nào không ghi khác thì dùng tỷ lệ của tổng.}
\khaibaokhainiem{diem-tuong-tac}{Điểm tương tác}
  {Điểm tương tác là điểm tổng hợp phản ánh mức người xem tương tác với nội dung của chiến dịch.}
  {}
  {\code{engagement_score}, thang điểm từ khoảng 2,5 tới 31; không phải tỷ lệ phần trăm và không cộng dồn giữa các chiến dịch.}
\khaibaokhainiem{chi-phi-chien-dich}{Chi phí chiến dịch}
  {Chi phí chiến dịch là toàn bộ số tiền đã chi cho một chiến dịch.}
  {acquisition\_cost × conversions}
  {dữ liệu không có cột tổng chi phí; \code{acquisition_cost} là chi phí trên một lượt chuyển đổi. Đơn vị tiền không được ghi.}
\khaibaokhainiem{roi}{ROI (lợi tức đầu tư)}
  {ROI phản ánh phần lợi nhuận thu về trên mỗi đơn vị chi phí marketing. ROI dương là có lãi, bằng 0 là hoà vốn, âm là lỗ.}
  {(doanh thu - chi phí) / chi phí}
  {cột \code{roi} của \code{campaigns} đã tính sẵn và làm tròn 2 chữ số. ROI của một nhóm tính từ tổng doanh thu và tổng chi phí, không lấy bình quân \code{roi}.}

% --- Dữ liệu thô
\khaibaokhainiem{du-lieu-tho}{Dữ liệu thô và bản sạch}
  {Dữ liệu thô là dữ liệu đúng như khi nhận từ nguồn, chưa kiểm tra hay sửa. Bản sạch là dữ liệu đã loại lỗi.}
  {}
  {\code{campaigns_raw} giữ nguyên ba file có nhiễu (chuỗi gốc, ô trống thành NULL, có dòng trùng); \code{campaigns} là bản sạch. Khoá dòng của bản thô là (\code{brand_id}, \code{row_index}).}
\khaibaokhainiem{nhat-ky-nhieu}{Nhật ký nhiễu}
  {Nhật ký nhiễu ghi lại từng lỗi đã được chèn vào dữ liệu thô để luyện làm sạch.}
  {}
  {\code{noise_log}: \code{noise_type} là chuỗi (\code{'NULL'}, \code{'FUNNEL_ERROR'}, \code{'DUPLICATE'}, \code{'INVALID_DATE'}, \code{'NEGATIVE'}, \code{'WHITESPACE'}). Nối với bản thô theo \code{brand_id} và \code{row_index}; với dòng trùng, \code{row_index} là vị trí của dòng gốc bị chép.}
\khaibaokhainiem{dong-trung}{Dòng trùng hoàn toàn}
  {Dòng trùng hoàn toàn là dòng giống một dòng khác ở mọi cột nghiệp vụ.}
  {}
  {so mọi cột của \code{campaigns_raw} trừ \code{raw_id} và \code{row_index} (hai cột vị trí luôn khác nhau).}

\begin{document}
{\LARGE\bfseries Gói câu hỏi SQL · Marketing · gói 3\par}
\medskip
{\color{cqMuted}Bộ dữ liệu \texttt{mkt} · Module 5 Joins (9 câu) · Module 6 Subqueries (9 câu) · case study Dễ và Trung bình · chuẩn \texttt{aio-sql-cauhoi} v0.4\par}
\bigskip
Lô soạn chờ kiểm và chờ giảng viên duyệt. Khung xanh nét đứt (\emph{Thiết kế ERD}) là ghi chú cho người duyệt. Các khung Truy vấn khởi đầu, Lời giải và Bài sai lấy từ \code{solution.sql}; học viên chỉ thấy truy vấn khởi đầu.

Mã danh mục dùng trong đề: thương hiệu \code{brand_id} 1 Nykaa, 2 Purplle, 3 Tira Beauty; loại chiến dịch \code{type_id} 1 Email, 2 Influencer, 3 Paid Ads, 4 SEO, 5 Social Media; kênh \code{channel_id} 1 Email, 2 Facebook, 3 Google, 4 Instagram, 5 WhatsApp, 6 YouTube.

\tableofcontents


% ################################################################ MODULE 5 · JOINS

% ================================================================ J1 · Viết · Cơ bản · A1 · sales
\begin{cauhoi}[ma=mkt-j1-tira-doanh-thu-cao-kem-loai, dulieu=mkt, domain=sales,
  hinhthuc=write, cap=basic, dang=A1, kynang={INNER JOIN, WHERE}, bay={noi-sai-khoa},
  tacgia={A}, trangthai=review, bokiemtra={public, hidden-1}, thutu=co, socot=vitri]
\tieude{Chiến dịch nào của Tira Beauty ngày 19/3/2025 có doanh thu từ 2.000.000?}
\boicanh[Chị Ngọc · Quản lý thương hiệu Tira Beauty]{Chị đang chọn vài chiến dịch khởi chạy ngày 19/3 có doanh thu tốt để làm mẫu cho đợt sau. Em lấy giúp chị danh sách chiến dịch có doanh thu từ 2 triệu (đơn vị tiền) trở lên, ghi rõ tên loại chiến dịch nhé.}
\begin{yeucau}
  \item Xét các chiến dịch của Tira Beauty (\code{brand_id = 3}) bắt đầu ngày 19/3/2025 có doanh thu (\code{revenue}) từ 2.000.000 trở lên.
  \item Nối \code{campaigns} với \code{campaign_types} theo \code{type_id} để lấy tên loại chiến dịch.
  \item Trả về \code{campaign_id}, \code{type_name}, \code{revenue}; sắp \code{revenue} giảm dần, bằng nhau thì \code{campaign_id} tăng dần (so theo chuỗi).
\end{yeucau}
\begin{erd}
  \erdbang{campaigns}{0,0}{campaign_id:PK='TI-CMP-18272', brand_id:FK*=3, type_id:FK*=4, start_date*=2025-03-18, revenue*=1964996}
  \erdbang{campaign_types}{2,0}{type_id:PK=4, type_name*='SEO'}
  \erdnoi[kieu=n-1, tren={\code{campaigns.type_id = campaign_types.type_id}}, duoi={nhiều chiến dịch, một loại}]{campaigns.type_id}{campaign_types.type_id}
\end{erd}
\begin{nghiepvu}{Chọn chiến dịch có doanh thu tốt của Tira Beauty để làm mẫu cho các đợt sau.}
  \khainiem{chien-dich}
  \luuy{Bảng \code{campaigns} chỉ lưu mã loại chiến dịch; tên loại nằm ở danh mục \code{campaign_types}. Doanh thu ghi bằng đơn vị tiền, dữ liệu không ghi loại tiền.}
\end{nghiepvu}
\goiy{\code{JOIN campaign_types t ON t.type_id = c.type_id}, rồi lọc thương hiệu, ngày và doanh thu ở \code{WHERE}.}
\ghichuthietke{Cơ bản (v4, hạ từ Dễ). Năm trục: bảng: lời giải 2 bảng, đúng mức Cơ bản của Joins trên lộ trình (data.js \code{j1-don-kem-vung}: nối 2 bảng, cấp basic); ERD vẽ đúng 2 bảng đó, có giá trị mẫu (chiến dịch ngày 18/3, không thuộc kết quả) và chữ trên đường nối, không bảng nhiễu; bước: 1 tầng (nối, lọc, sắp), không gom; mơ hồ: yêu cầu ghi đủ cột, giá trị lọc và khoá nối; bẫy: 0; nghiệp vụ: một khái niệm, không công thức. Đạt Cơ bản 5/5. Bản Dễ cũ (số chiến dịch và doanh thu theo thương hiệu và loại quý 2, nối 3 bảng, gom) đã bỏ. Kết quả public 7 dòng, hidden-1 6 dòng (bỏ TI-CMP-24167). Wrong: \code{noi_sai_khoa} (nối loại theo \code{brand_id}, mọi dòng mang tên loại mã 3), \code{tra_ma_loai} (không nối danh mục, trả mã loại); cả hai rớt public. Ô: A1 · sales · basic · write, không trùng trong gói.}
\end{cauhoi}

% ================================================================ J2 · Viết · Dễ · A3 · customer
\begin{cauhoi}[ma=mkt-j2-doi-tuong-va-phan-khuc, dulieu=mkt, domain=customer,
  hinhthuc=write, cap=easy, dang=A3, kynang={JOIN hai lần một bảng, bí danh}, bay={noi-mot-lan-cho-hai-vai},
  tacgia={A}, trangthai=review, bokiemtra={public, hidden-1}, thutu=co, socot=vitri]
\tieude{Đối tượng mục tiêu và phân khúc khách hàng khớp nhau đến đâu?}
\boicanh[Chị Hoa · Chuyên viên phân tích marketing]{Chị đang xem chiến dịch nhắm tới một nhóm thì khách thực tế mua hàng thường thuộc nhóm nào. Chị cần số chiến dịch của từng cặp đối tượng mục tiêu và phân khúc khách hàng trong 6 tháng cuối năm 2024, có ghi tên nhóm cho dễ đọc.}
\begin{yeucau}
  \item Xét các chiến dịch bắt đầu từ ngày 1/7/2024 tới hết ngày 31/12/2024.
  \item Đếm số chiến dịch theo từng cặp đối tượng mục tiêu và phân khúc khách hàng (ghi bằng tên nhóm).
  \item Trả về \code{doi_tuong}, \code{phan_khuc} (tên nhóm trong \code{audiences}), \code{so_chien_dich}.
  \item Sắp theo \code{doi_tuong} rồi \code{phan_khuc}, tăng dần.
\end{yeucau}
\begin{erd}
  \erdbang{campaigns}{0,0}{campaign_id:PK, target_audience_id:FK*, segment_id:FK*, start_date*}
  \erdbang{audiences}{2,0}{audience_id:PK, audience_name*}
  \erdnoi[kieu=n-1, tren={vai 1: đối tượng mục tiêu}]{campaigns.target_audience_id}{audiences.audience_id}
  \erdnoi[kieu=n-1, tren={vai 2: phân khúc khách hàng}]{campaigns.segment_id}{audiences.audience_id}
\end{erd}
\begin{nghiepvu}{Đánh giá chiến dịch nhắm một nhóm có thực sự đưa về khách thuộc đúng nhóm đó không, để chỉnh cách chọn đối tượng.}
  \khainiem{doi-tuong-va-phan-khuc}
  \luuy{Một danh mục dùng cho hai vai: mỗi vai cần một lần nối riêng với một bí danh riêng. Một lần nối có cả hai điều kiện chỉ giữ chiến dịch có hai cột bằng nhau.}
\end{nghiepvu}
\goiy{Muốn có hai tên khác nhau trên cùng một dòng thì bảng \code{audiences} phải xuất hiện mấy lần trong \code{FROM}?}
\goiy{\code{JOIN audiences dt ON dt.audience_id = c.target_audience_id JOIN audiences pk ON pk.audience_id = c.segment_id}.}
\ghichuthietke{Dễ: 2 bảng, hai đường nối có chữ ghi vai. Kết quả 25 dòng. Bẫy ở public: một lần nối với cả hai điều kiện ra 5 dòng; trả mã phân khúc thay cho tên. hidden-1 bỏ chiến dịch đuôi 7.}
\end{cauhoi}

% ================================================================ J3 · Viết · Cơ bản · A1 · marketing
\begin{cauhoi}[ma=mkt-j3-influencer-nykaa-tren-youtube, dulieu=mkt, domain=marketing,
  hinhthuc=write, cap=basic, dang=A1, kynang={INNER JOIN, bảng cầu nối}, bay={chi-lay-kenh-dau},
  tacgia={A}, trangthai=review, bokiemtra={public, hidden-1}, thutu=co, socot=vitri]
\tieude{Chiến dịch Influencer nào của Nykaa ngày 15/5/2025 có chạy trên YouTube?}
\boicanh[Chị Trang · Quản lý thương hiệu Nykaa]{Chị chuẩn bị buổi làm việc với đối tác YouTube về các chiến dịch người ảnh hưởng. Em lấy giúp chị các chiến dịch Influencer khởi chạy ngày 15/5 có dùng kênh YouTube, kèm lượt hiển thị.}
\begin{yeucau}
  \item Xét các chiến dịch của Nykaa (\code{brand_id = 1}), loại Influencer (marketing người ảnh hưởng, \code{type_id = 2}), bắt đầu ngày 15/5/2025.
  \item Nối \code{campaigns} với bảng cầu nối \code{campaign_channels} theo \code{campaign_id}; giữ các chiến dịch có kênh YouTube (\code{channel_id = 6}), dù YouTube đứng ở vị trí nào trong danh sách kênh.
  \item Trả về \code{campaign_id}, \code{impressions}; sắp \code{impressions} giảm dần, bằng nhau thì \code{campaign_id} tăng dần.
\end{yeucau}
\begin{erd}
  \erdbang{campaigns}{0,0}{campaign_id:PK='NY-CMP-16407', brand_id:FK*=1, type_id:FK*=2, start_date*=2025-05-14, impressions*=79221}
  \erdbang{campaign_channels}{2,0}{campaign_id:FK='NY-CMP-16407', channel_id:FK*=6, channel_pos=3}
  \erdnoi[kieu=1-n, tren={\code{campaigns.campaign_id = campaign_channels.campaign_id}}, duoi={một chiến dịch, 1 tới 3 kênh}]{campaigns.campaign_id}{campaign_channels.campaign_id}
\end{erd}
\begin{nghiepvu}{Chuẩn bị số liệu làm việc với đối tác YouTube về các chiến dịch người ảnh hưởng của Nykaa.}
  \khainiem{kenh-da-tri}
  \luuy{Mỗi dòng của \code{campaign_channels} là một cặp chiến dịch, kênh. \code{channel_pos} chỉ là thứ tự kênh trong chuỗi gốc; YouTube đứng thứ hai hay thứ ba thì vẫn là kênh của chiến dịch.}
\end{nghiepvu}
\goiy{\code{JOIN campaign_channels cc ON cc.campaign_id = c.campaign_id}, rồi lọc \code{cc.channel_id = 6} cùng thương hiệu, loại và ngày.}
\ghichuthietke{Cơ bản (v4, hạ từ Trung bình). Năm trục: bảng: lời giải 2 bảng (Cơ bản của Joins trên lộ trình, data.js \code{j1-don-kem-vung}); ERD 2 bảng có giá trị mẫu (chiến dịch ngày 14/5 có YouTube ở vị trí 3, không thuộc kết quả) và chữ trên đường nối, bỏ bảng \code{channels}; bước: 1 tầng (nối, lọc, sắp), không gom; mơ hồ: yêu cầu ghi đủ cột, giá trị lọc, khoá nối và cách hiểu vị trí kênh; bẫy: 0 (ý vị trí kênh đã nói trong yêu cầu); nghiệp vụ: một khái niệm. Đạt Cơ bản 5/5. Bản Trung bình cũ (đếm chiến dịch và cộng lượt hiển thị theo kênh, 3 bảng) đã bỏ. Kết quả public 8 dòng, hidden-1 7 dòng (bỏ NY-CMP-17647). Wrong: \code{chi_lay_kenh_dau} (thêm \code{channel_pos = 1}, public còn 4 dòng), \code{noi_khong_loc_kenh} (mỗi chiến dịch Influencer một dòng cho mỗi kênh); cả hai rớt public. Ô: A1 · marketing · basic · write, không trùng trong gói (cùng ô với \code{mkt-r1-influencer-purplle-thang-1} ở gói 1, khác module).}
\end{cauhoi}

% ================================================================ J4 · Viết · Trung bình · B4 · operations
\begin{cauhoi}[ma=mkt-j4-dang-chay-tuan-dau-thang-2, dulieu=mkt, domain=operations,
  hinhthuc=write, cap=medium, dang=B4, kynang={nối theo khoảng ngày, DATE\_ADD, GROUP BY}, bay={lech-bien-ngay-cuoi, noi-bang},
  tacgia={A}, trangthai=review, bokiemtra={public, hidden-1}, thutu=co, socot=vitri]
\tieude{Mỗi ngày tuần 3--9/2/2025 có bao nhiêu chiến dịch đang chạy?}
\boicanh[Chị Hạnh · Trưởng phòng marketing số]{Tuần trước Valentine đội vận hành quảng cáo kêu quá tải. Chị muốn biết mỗi ngày từ thứ Hai 3/2 tới Chủ nhật 9/2 có bao nhiêu chiến dịch đang chạy cùng lúc, tính cả ba thương hiệu.}
\begin{yeucau}
  \item Đếm số chiến dịch đang chạy của từng ngày trong lịch ngày, từ 3/2/2025 tới hết 9/2/2025.
  \item Chiến dịch đang chạy vào ngày d khi d nằm trong khoảng từ ngày bắt đầu tới ngày cuối, tính cả hai đầu; ngày cuối = \code{start_date} + \code{duration_days} - 1 ngày.
  \item Trả về \code{cal_date}, \code{so_dang_chay}; sắp \code{cal_date} tăng dần.
\end{yeucau}
\begin{erd}
  \erdbang{calendar}{0,0}{cal_date:PK, day_name, week_start}
  \erdbang{campaigns}{2,0}{campaign_id:PK, start_date:FK*, duration_days*, type_id:FK}
  \erdbang{campaign_types}{2,1}{type_id:PK, type_name}
  \erdnoi[kieu=1-n]{calendar.cal_date}{campaigns.start_date}
  \erdnoi[kieu=n-1]{campaigns.type_id}{campaign_types.type_id}
\end{erd}
\begin{nghiepvu}{Biết khối lượng chiến dịch phải theo dõi mỗi ngày để bố trí người trực vận hành quảng cáo.}
  \quytrinh{Vòng đời chiến dịch}{lên kế hoạch, khởi chạy, đang chạy, kết thúc, đánh giá}{}
  \khainiem{chien-dich-dang-chay}
  \luuy{Đường nối trong ERD là quan hệ khoá (ngày bắt đầu là một ngày của lịch). Đếm chiến dịch đang chạy cần so ngày của lịch với cả khoảng chạy, từ ngày bắt đầu tới ngày cuối. Chiến dịch 5 ngày bắt đầu ngày 1 thì chạy tới hết ngày 5, không tới ngày 6.}
\end{nghiepvu}
\goiy{Điều kiện nối không nhất thiết là dấu bằng: một ngày lịch khớp với mọi chiến dịch có khoảng chạy chứa ngày đó.}
\goiy{\code{ON c.start_date <= k.cal_date AND k.cal_date <= DATE_ADD(c.start_date, INTERVAL c.duration_days - 1 DAY)}.}
\ghichuthietke{Trung bình: 3 bảng (\code{campaign_types} là bảng nhiễu, theo CHOT mục 3), 2 bảng cần dùng nhưng nối theo khoảng ngày, không theo khoá, đường nối 1-n là quan hệ khoá gây nhiễu có chủ đích. Bẫy ở public: quên trừ 1 ngày (thừa chiến dịch kết thúc hôm trước); nối bằng \code{=} chỉ đếm chiến dịch khởi chạy trong ngày. hidden-1 bỏ chiến dịch đuôi 7 nên chép số sẽ rớt.}
\end{cauhoi}

% ================================================================ J5 · Viết · Trung bình · A1 · operations
\begin{cauhoi}[ma=mkt-j5-ngay-khong-co-chien-dich, dulieu=mkt, domain=operations,
  hinhthuc=write, cap=medium, dang=A1, kynang={LEFT JOIN, điều kiện trong ON, IS NULL}, bay={noi-trong, dieu-kien-o-where},
  tacgia={A}, trangthai=review, bokiemtra={public, hidden-1}, thutu=co, socot=vitri]
\tieude{Ngày nào tháng 5/2025 Tira Beauty không mở Paid Ads cho Premium Shoppers?}
\boicanh[Anh Tùng · Chuyên viên phân tích marketing]{Tira Beauty muốn ngày nào cũng có quảng cáo trả phí (Paid Ads) tiếng Hindi nhắm khách mua sắm cao cấp (Premium Shoppers). Em tìm giúp anh những ngày tháng 5 vừa rồi không có chiến dịch nào như vậy khởi chạy.}
\begin{yeucau}
  \item Xét mọi ngày từ 1/5/2025 tới hết 31/5/2025 trong lịch ngày \code{calendar}.
  \item Liệt kê các ngày không có chiến dịch nào của Tira Beauty (\code{brand_id = 3}), loại Paid Ads (\code{type_id = 3}), ngôn ngữ Hindi (\code{language_id = 3}), đối tượng mục tiêu Premium Shoppers (\code{target_audience_id = 2}) bắt đầu vào ngày đó.
  \item Trả về \code{cal_date}, \code{day_name}; sắp \code{cal_date} tăng dần.
\end{yeucau}
\begin{erd}
  \erdbang{calendar}{0,0}{cal_date:PK, day_name*, is_weekend}
  \erdbang{campaigns}{2,0}{campaign_id:PK, start_date:FK*, brand_id:FK*, type_id*, language_id*, target_audience_id*}
  \erdbang{brands}{2,1}{brand_id:PK, brand_name}
  \erdnoi[kieu=1-n]{calendar.cal_date}{campaigns.start_date}
  \erdnoi[kieu=n-1]{campaigns.brand_id}{brands.brand_id}
\end{erd}
\begin{nghiepvu}{Tìm ngày bị bỏ trống trong lịch quảng cáo trả phí của một nhóm khách mục tiêu để lên lịch bù.}
  \khainiem{chien-dich}
  \khainiem{doi-tuong-va-phan-khuc}
  \luuy{Ngày không có chiến dịch không có dòng nào trong \code{campaigns}, nên phải đi từ lịch ngày. Sau nối ngoài trái, ngày không khớp với chiến dịch nào vẫn giữ một dòng với các cột của \code{campaigns} là NULL. Tháng 5/2025 là tháng đủ dữ liệu.}
\end{nghiepvu}
\goiy{Nối trong (\code{JOIN}) có giữ lại ngày không khớp với chiến dịch nào không? Nếu đặt điều kiện thương hiệu, loại ở WHERE thì dòng NULL của những ngày đó còn không?}
\goiy{\code{calendar k LEFT JOIN campaigns c ON c.start_date = k.cal_date AND c.brand_id = 3 AND ...} rồi \code{WHERE c.campaign_id IS NULL}.}
\ghichuthietke{Trung bình: 3 bảng (\code{brands} là bảng nhiễu, theo CHOT mục 3), ký hiệu 1-n; độ khó ở nối ngoài có điều kiện trong ON rồi lọc dòng không khớp. Sửa sau kiểm B: bản cũ (mọi chiến dịch, cả lịch) ra 6 ngày 25--30/6 là ngày hết dữ liệu, giống nhau ở hai biến thể và lộ qua khai báo dữ liệu. Bản mới thu hẹp nhóm, chọn tháng đủ: public 3 ngày, hidden-1 5 ngày (bỏ chiến dịch đuôi 7 làm lộ thêm ngày trống). Bẫy ở public: nối trong ra 0 dòng; điều kiện ở WHERE ra 0 dòng; bỏ điều kiện ngôn ngữ thì không còn ngày trống. Phân biệt với q6 (NOT EXISTS, Purplle SEO) và j9 (đếm theo ngày).}
\end{cauhoi}

% ================================================================ J6 · Viết · Khó · E2 · marketing
\begin{cauhoi}[ma=mkt-j6-email-lech-kenh-nua-dau-2025, dulieu=mkt, domain=marketing,
  hinhthuc=write, cap=hard, dang=E2, kynang={LEFT JOIN có điều kiện trong ON, đếm có điều kiện, nối 4 bảng}, bay={nham-loai-voi-kenh, fan-out, dieu-kien-kenh-o-where, sai-mau-so},
  tacgia={A}, trangthai=review, bokiemtra={public, hidden-1}, thutu=co, socot=vitri]
\tieude{Bao nhiêu chiến dịch loại Email nửa đầu 2025 không chạy kênh Email?}
\boicanh[Chị Thư · Chuyên viên email marketing]{Giám đốc marketing hỏi vì sao báo cáo theo kênh Email có ít chiến dịch hơn hẳn báo cáo theo loại Email. Em tính giúp chị, cho nửa đầu 2025 và theo từng thương hiệu, có bao nhiêu chiến dịch email lệch kênh, chiếm tỷ lệ bao nhiêu và mang về bao nhiêu doanh thu.}
\begin{yeucau}
  \item Xét các chiến dịch loại Email bắt đầu từ 1/1/2025 tới hết 30/6/2025.
  \item Theo từng thương hiệu, tính số chiến dịch loại Email, số chiến dịch email lệch kênh (định nghĩa ở tab Nghiệp vụ), tỷ lệ lệch kênh và doanh thu của nhóm lệch kênh.
  \item Trả về \code{brand_name}, \code{so_chien_dich_email}, \code{so_lech_kenh}, \code{ty_le_lech_pct} (phần trăm trên số chiến dịch loại Email của thương hiệu, làm tròn 2 chữ số), \code{doanh_thu_lech} (tổng \code{revenue} của nhóm lệch kênh); sắp \code{brand_name} tăng dần.
\end{yeucau}
\begin{erd}
  \erdbang{brands}{0,0}{brand_id:PK}
  \erdbang{campaigns}{1,0}{campaign_id:PK, brand_id:FK, type_id:FK, target_audience_id:FK}
  \erdbang{campaign_types}{2,0}{type_id:PK}
  \erdbang{campaign_channels}{1,1}{campaign_id:FK, channel_id:FK}
  \erdbang{channels}{2,1}{channel_id:PK}
  \erdbang{audiences}{0,1}{audience_id:PK}
  \erdkhoa{\code{campaigns} nối \code{brands} theo \code{brand_id} và \code{campaign_types} theo \code{type_id}; \code{campaign_channels} nối \code{campaigns} theo \code{campaign_id} và \code{channels} theo \code{channel_id}. \code{audiences} không cần; \code{channels} chỉ để tra mã kênh Email.}
\end{erd}
\begin{nghiepvu}{Quyết định có thống nhất lại cách gán loại Email cho chiến dịch hay không, trước khi chia ngân sách email nửa cuối năm.}
  \khainiem{kenh-da-tri}
  \luuy{Chiến dịch email lệch kênh là chiến dịch thuộc loại Email nhưng trong danh sách kênh của chiến dịch không có kênh Email. Chiến dịch có kênh Email kèm Facebook hay YouTube không tính là lệch kênh. Loại do người lập kế hoạch gán cho cả chiến dịch; kênh là nơi chiến dịch thực sự xuất hiện. Doanh thu ghi cho cả chiến dịch, nên mỗi chiến dịch chỉ được cộng doanh thu một lần dù chạy mấy kênh.}
\end{nghiepvu}
\goiy{Ở mỗi chiến dịch loại Email, chỉ cần biết bảng cầu nối có dòng kênh Email của chiến dịch đó hay không. Nối với mọi dòng kênh rồi mới gom thì doanh thu bị cộng theo số kênh.}
\goiy{\code{LEFT JOIN campaign_channels e ON e.campaign_id = c.campaign_id AND e.channel_id = 1}; chiến dịch lệch kênh là dòng có \code{e.campaign_id IS NULL}.}
\ghichuthietke{Khó (v4, nâng từ Trung bình). Năm trục: bảng: lời giải 4 bảng (\code{brands}, \code{campaigns}, \code{campaign_types}, \code{campaign_channels}); ERD 6 bảng chỉ khoá, 2 bảng nhiễu (\code{channels} chỉ để tra mã, \code{audiences}), không vẽ đường nối, cách nối ở \code{erdkhoa}; bước: 4 tầng (đánh dấu từng chiến dịch có kênh Email hay không, gom theo thương hiệu, tỷ lệ, doanh thu có điều kiện); mơ hồ: chỉ số ``lệch kênh'' định nghĩa ở tab Nghiệp vụ, lời nhắn nói bằng ngôn ngữ báo cáo; bẫy: 4 (lọc kênh khác Email ở WHERE, doanh thu nhân theo số kênh khi nối thẳng cầu nối, đảo vai loại và kênh, mẫu số là mọi chiến dịch); nghiệp vụ: chỉ số nhiều điều kiện (loại, kênh, tỷ lệ, doanh thu). Đạt Khó 4/5; chưa có bẫy chỉ lộ ở hidden-1 vì chiến dịch thêm ngày 1/7/2025 của hidden-1 thuộc loại Social Media. Bản Trung bình cũ (cả kỳ, chỉ hai cột đếm) đã bỏ. Kết quả public khác hidden-1 (bỏ chiến dịch đuôi 7). Wrong: \code{kenh_khac_email_o_where}, \code{doanh_thu_nhan_dong} (đếm đúng nhờ \code{COUNT(DISTINCT)} nhưng doanh thu lấy tổng mọi dòng nối trừ dòng kênh Email), \code{kenh_email_khong_phai_loai_email}, \code{mau_so_moi_chien_dich}; cả bốn rớt public. Cách viết \code{NOT EXISTS}, \code{NOT IN}, bảng dẫn xuất \code{MAX(channel_id = 1)} cho cùng kết quả, đều dưới 2 giây. Ô: E2 · marketing · hard · write, không trùng trong gói.}
\end{cauhoi}

% ================================================================ J7 · Sửa lỗi · Dễ · A3 · sales
\begin{cauhoi}[ma=mkt-j7-doanh-thu-instagram-facebook-gap-doi, dulieu=mkt, domain=sales,
  hinhthuc=debug, cap=easy, dang=A3, kynang={bảng cầu nối, nhân dòng, bảng dẫn xuất}, bay={fan-out},
  tacgia={A}, trangthai=review, bokiemtra={public, hidden-1}, thutu=co, socot=vitri]
\tieude{Doanh thu chiến dịch dùng Instagram hoặc Facebook bị cộng hai lần}
\boicanh[Anh Vinh · Trưởng nhóm quảng cáo trả phí]{Báo cáo quý 1 của Tira Beauty ghi doanh thu các chiến dịch có Instagram hoặc Facebook cao hơn thực tế. Em tìm xem số sai ở đâu rồi báo anh.}
\begin{yeucau}
  \item Truy vấn khởi đầu đếm số chiến dịch của Tira Beauty (\code{brand_id = 3}) bắt đầu từ 1/1/2025 tới hết 31/3/2025 có dùng kênh Facebook (\code{channel_id = 2}) hoặc Instagram (\code{channel_id = 4}).
  \item Truy vấn cũng cộng doanh thu của các chiến dịch đó. Mỗi chiến dịch chỉ được tính một lần.
  \item Sửa truy vấn cho đúng yêu cầu; giữ hai cột \code{so_chien_dich}, \code{doanh_thu}.
\end{yeucau}
\begin{erd}
  \erdbang{campaigns}{0,0}{campaign_id:PK, brand_id*, start_date*, revenue*}
  \erdbang{campaign_channels}{2,0}{campaign_id:FK, channel_id:FK*}
  \erdnoi[kieu=1-n, tren={một chiến dịch, 1 tới 3 kênh}, canhbao]{campaigns.campaign_id}{campaign_channels.campaign_id}
\end{erd}
\begin{nghiepvu}{Báo cáo đúng doanh thu của nhóm chiến dịch có mặt trên hai mạng xã hội lớn, trước khi quyết định ngân sách quý 2.}
  \khainiem{kenh-da-tri}
  \luuy{Chiến dịch chạy cả Facebook lẫn Instagram có hai dòng khớp điều kiện trong bảng cầu nối. Doanh thu ghi cho cả chiến dịch, nên nối thẳng rồi cộng sẽ cộng doanh thu của chiến dịch đó hai lần.}
\end{nghiepvu}
\goiy{Sau khi nối, một chiến dịch có thể xuất hiện mấy dòng? So \code{COUNT(*)} với \code{COUNT(DISTINCT c.campaign_id)}.}
\goiy{Lấy danh sách \code{campaign_id} khác nhau có kênh 2 hoặc 4 trong một bảng dẫn xuất (\code{SELECT DISTINCT}) rồi mới nối với \code{campaigns}.}
\ghichuthietke{Dễ: 2 bảng, chữ trên đường nối tô cảnh báo quan hệ một chiến dịch nhiều kênh. Bẫy ở public: chỉ đổi \code{COUNT(DISTINCT)} mà vẫn \code{SUM} trên dòng đã nhân; \code{SUM(DISTINCT revenue)} bỏ cả chiến dịch khác nhau có cùng doanh thu. Lời giải cũng có thể viết bằng \code{IN} hoặc \code{EXISTS} (Module 6).}
\end{cauhoi}

% ================================================================ J8 · Sửa lỗi · Dễ · A3 · customer
\begin{cauhoi}[ma=mkt-j8-phan-khuc-lay-nham-doi-tuong, dulieu=mkt, domain=customer,
  hinhthuc=debug, cap=easy, dang=A3, kynang={khoá ngoại hai vai, JOIN}, bay={nham-vai-khoa-ngoai},
  tacgia={A}, trangthai=review, bokiemtra={public, hidden-1}, thutu=co, socot=vitri]
\tieude{Doanh thu theo phân khúc giống hệt doanh thu theo đối tượng}
\boicanh[Chị Ngân · Quản lý thương hiệu Purplle]{Chị có hai báo cáo nửa đầu 2025 của Purplle: một theo đối tượng mục tiêu, một theo phân khúc khách hàng. Hai bảng ra giống hệt nhau từng con số, chị nghĩ báo cáo phân khúc bị lập sai.}
\begin{yeucau}
  \item Báo cáo cần số chiến dịch và tổng doanh thu theo \textbf{phân khúc khách hàng} của Purplle (\code{brand_id = 2}), xét các chiến dịch bắt đầu từ 1/1/2025 tới hết 30/6/2025.
  \item Sửa truy vấn cho đúng yêu cầu; giữ các cột \code{phan_khuc}, \code{so_chien_dich}, \code{doanh_thu} và thứ tự sắp theo \code{phan_khuc}.
\end{yeucau}
\begin{erd}
  \erdbang{campaigns}{0,0}{campaign_id:PK, brand_id, target_audience_id:FK, segment_id:FK, start_date, revenue}
  \erdbang{audiences}{2,0}{audience_id:PK, audience_name}
  \erdnoi[kieu=n-1]{campaigns.target_audience_id}{audiences.audience_id}
  \erdnoi[kieu=n-1]{campaigns.segment_id}{audiences.audience_id}
\end{erd}
\begin{nghiepvu}{Biết nhóm khách nào thực sự mang về doanh thu cho Purplle, tách khỏi nhóm mà chiến dịch nhắm tới.}
  \khainiem{doi-tuong-va-phan-khuc}
  \luuy{Hai khoá ngoại cùng trỏ về một danh mục. Nối nhầm cột vẫn chạy được, vẫn ra đủ 5 nhóm với tên hợp lệ, nên lỗi chỉ lộ khi so với báo cáo kia.}
\end{nghiepvu}
\goiy{Trong \code{campaigns} cột nào ghi phân khúc khách hàng? Truy vấn đang nối \code{audiences} qua cột nào?}
\ghichuthietke{Dễ (hạ từ Trung bình sau kiểm C: lỗi một khoá nối, lời nhắn đã chỉ triệu chứng): 2 bảng, hai đường nối n-1 không ghi chữ, không tô cảnh báo. Bẫy ở public: nối cả hai điều kiện chỉ giữ chiến dịch có hai cột bằng nhau; sửa khoá nhưng bỏ lọc thương hiệu. hidden-1 bỏ chiến dịch đuôi 7.}
\end{cauhoi}

% ================================================================ J9 · Sửa lỗi · Khó · B1 · operations
\begin{cauhoi}[ma=mkt-j9-lich-paid-ads-tira-chua-dung-duoc, dulieu=mkt, domain=operations,
  hinhthuc=debug, cap=hard, dang=B1, kynang={LEFT JOIN, điều kiện trong ON, COUNT cột}, bay={dieu-kien-o-where, count-sao, thieu-can-tren},
  tacgia={A}, trangthai=review, bokiemtra={public, hidden-1}, thutu=co, socot=vitri]
\tieude{Lịch khởi chạy Paid Ads tháng 6/2025 của Tira Beauty chưa dùng được}
\boicanh[Anh Vinh · Trưởng nhóm quảng cáo trả phí]{Anh cần bảng số chiến dịch quảng cáo trả phí (Paid Ads) của Tira Beauty khởi chạy mỗi ngày trong tháng 6/2025, đủ 30 ngày, ngày nào không có thì ghi 0, để vẽ biểu đồ và dùng lại cho các tháng sau. Bảng bạn thực tập lập còn thiếu ngày. Em sửa giúp anh.}
\begin{yeucau}
  \item Kết quả cần đúng một dòng cho mỗi ngày từ 1/6/2025 tới hết 30/6/2025 trong lịch ngày, kèm số chiến dịch Paid Ads (\code{type_id = 3}) của Tira Beauty (\code{brand_id = 3}) bắt đầu ngày đó; ngày không có chiến dịch ghi 0.
  \item Truy vấn khởi đầu có nhiều chỗ sai chồng nhau; sửa một chỗ chưa đủ.
  \item Giữ cột \code{cal_date}, \code{so_chien_dich}, sắp \code{cal_date} tăng dần.
\end{yeucau}
\begin{erd}
  \erdbang{calendar}{0,0}{cal_date:PK}
  \erdbang{campaigns}{1,0}{campaign_id:PK, brand_id:FK, type_id:FK, start_date:FK}
  \erdbang{brands}{2,0}{brand_id:PK}
  \erdbang{campaign_types}{1,1}{type_id:PK}
  \erdbang{campaign_channels}{2,1}{campaign_id:FK, channel_id:FK}
  \erdkhoa{\code{campaigns.start_date} là một ngày của \code{calendar}; \code{brand_id}, \code{type_id} trỏ tới \code{brands}, \code{campaign_types}. Lời giải chỉ cần \code{calendar} và \code{campaigns}; ba bảng còn lại không cần.}
\end{erd}
\begin{nghiepvu}{Có lịch khởi chạy đủ ngày để vẽ biểu đồ nhịp mở quảng cáo trả phí trong tháng.}
  \khainiem{ky-va-thang-du}
  \luuy{Tháng 6/2025 chưa đủ dữ liệu: các ngày sau ngày dữ liệu cuối vẫn phải có mặt trên biểu đồ với số 0. Lịch \code{calendar} sẽ được nối thêm các ngày mới khi có dữ liệu tháng sau.}
\end{nghiepvu}
\goiy{Một chỗ sai làm mất ngày không có chiến dịch; khi ngày đó hiện lại, cách đếm hiện tại cho nó số 1 thay vì 0. Phạm vi ngày của lịch cũng cần xem lại.}
\goiy{Đưa điều kiện thương hiệu, loại vào \code{ON}, đếm \code{COUNT(c.campaign_id)} và thêm cận trên \code{k.cal_date < '2025-07-01'}.}
\ghichuthietke{Khó (v4, nâng từ Trung bình). Năm trục: bảng: lời giải 2 bảng (mẫu d4 Khó dùng 2 bảng); ERD 5 bảng chỉ khoá, 3 bảng nhiễu cho đủ 5 bảng (QUYET-DINH-CHUNG 1.3), không vẽ đường nối, không tô cảnh báo, cách nối ở \code{erdkhoa}; sửa lỗi: ba lỗi chồng nhau (điều kiện của bảng bên phải đặt ở WHERE làm mất ngày trống; \code{COUNT(*)} cho ngày trống ra 1, chỉ lộ sau khi sửa lỗi thứ nhất; lọc lịch thiếu cận trên, chỉ lộ ở hidden-1 có thêm ngày 2025-07-01); bẫy: 3, một bẫy chỉ lộ ở hidden-1; mơ hồ: yêu cầu không chỉ chỗ sai, ERD trung tính; nghiệp vụ: lịch đủ ngày cho tháng chưa đủ dữ liệu. Đạt Khó 3/5 cộng trục sửa lỗi. Starter ra 24 dòng (mất 6 ngày cuối tháng). Wrong: \code{chi_dua_dieu_kien_vao_on} (rớt public: ngày trống ra 1; hidden-1 thêm dòng 1/7), \code{chi_sua_cot_dem} (rớt public), \code{sua_hai_loi_thieu_can_tren} (chỉ rớt hidden-1, như thiết kế). Bản Trung bình cũ có một lỗi trong starter. Ô: B1 · operations · hard · debug, không trùng trong gói.}
\end{cauhoi}


% ################################################################ MODULE 6 · SUBQUERIES

% ================================================================ Q1 · Viết · Dễ · C3 · finance
\begin{cauhoi}[ma=mkt-q1-roi-tren-binh-quan-tira, dulieu=mkt, domain=finance,
  hinhthuc=write, cap=easy, dang=C3, kynang={truy vấn con một giá trị}, bay={muc-chung-sai-pham-vi, avg-cua-ty-le},
  tacgia={A}, trangthai=review, bokiemtra={public, hidden-1}, thutu=co, socot=vitri]
\tieude{Bao nhiêu chiến dịch Tira Beauty tháng 4/2025 có ROI trên mức chung?}
\boicanh[Chị Mai · Kế toán quản trị]{Chị đang chấm điểm hiệu quả các chiến dịch tháng 4 của Tira Beauty. Em đếm giúp chị có bao nhiêu chiến dịch có ROI cao hơn ROI chung của cả thương hiệu trong tháng đó.}
\begin{yeucau}
  \item Xét các chiến dịch của Tira Beauty (\code{brand_id = 3}) bắt đầu từ 1/4/2025 tới hết 30/4/2025.
  \item Mức chung là ROI (lợi tức đầu tư) của cả nhóm chiến dịch này, tính theo tỷ lệ của tổng: (tổng doanh thu - tổng chi phí) / tổng chi phí, với chi phí mỗi chiến dịch = \code{acquisition_cost} × \code{conversions}.
  \item Đếm số chiến dịch có cột \code{roi} lớn hơn mức chung. Trả về một dòng một cột \code{so_chien_dich}.
\end{yeucau}
\begin{erd}
  \erdbang{campaigns}{0,0}{campaign_id:PK, brand_id*, start_date*, revenue*, acquisition_cost*, conversions*, roi*}
  \erdbang{brands}{2,0}{brand_id:PK, brand_name}
  \erdnoi[kieu=n-1, tren={\code{campaigns.brand_id = brands.brand_id}}]{campaigns.brand_id}{brands.brand_id}
\end{erd}
\begin{nghiepvu}{Tách nhóm chiến dịch hiệu quả hơn mức chung của thương hiệu để xem lại cách phân bổ chi phí tháng sau.}
  \khainiem{roi}
  \khainiem{chi-phi-chien-dich}
  \khainiem{binh-quan-gia-quyen}
  \luuy{Mức chung phải tính trên đúng nhóm đang xét (một thương hiệu, một tháng). Bình quân cột \code{roi} cho mức khác hẳn ROI của tổng, vì chiến dịch chi phí nhỏ thường có ROI rất cao, làm bình quân giản đơn cao hơn.}
\end{nghiepvu}
\goiy{Mức chung là một con số. Có thể tính con số này trong một truy vấn con đặt ngay sau dấu \code{>}.}
\goiy{Truy vấn con lặp lại đúng điều kiện thương hiệu và tháng của truy vấn ngoài, rồi tính \code{(SUM(revenue) - SUM(acquisition_cost * conversions)) / SUM(acquisition_cost * conversions)}.}
\ghichuthietke{Dễ: 2 bảng (\code{brands} chỉ để tra mã), chữ trên đường nối. Bẫy ở public: mức chung lấy cả kỳ và cả ba thương hiệu; mức chung là \code{AVG(roi)} (bình quân giản đơn, cao hơn). Ba cách cho ba số khác nhau. hidden-1 bỏ chiến dịch đuôi 7.}
\end{cauhoi}

% ================================================================ Q2 · Viết · Dễ · C1 · sales
\begin{cauhoi}[ma=mkt-q2-doanh-thu-cao-nhat-moi-thuong-hieu, dulieu=mkt, domain=sales,
  hinhthuc=write, cap=easy, dang=C1, kynang={truy vấn con tương quan, MAX}, bay={khong-tuong-quan, gom-nhom-lay-ma-sai},
  tacgia={A}, trangthai=review, bokiemtra={public, hidden-1}, thutu=co, socot=vitri]
\tieude{Chiến dịch nào doanh thu cao nhất của mỗi thương hiệu ngày 5/3/2025?}
\boicanh[Anh Tùng · Chuyên viên phân tích marketing]{Anh đang viết bài chia sẻ nội bộ về đợt chiến dịch chuẩn bị cho dịp 8/3, cần chiến dịch có doanh thu cao nhất của từng thương hiệu trong số các chiến dịch khởi chạy ngày 5/3.}
\begin{yeucau}
  \item Xét các chiến dịch bắt đầu ngày 5/3/2025.
  \item Lấy chiến dịch có doanh thu (\code{revenue}) cao nhất của từng thương hiệu; nếu nhiều chiến dịch cùng mức cao nhất thì lấy hết.
  \item Trả về \code{brand_id}, \code{campaign_id}, \code{revenue}; sắp \code{brand_id} rồi \code{campaign_id} tăng dần.
\end{yeucau}
\begin{erd}
  \erdbang{campaigns}{0,0}{campaign_id:PK, brand_id:FK*, start_date*, revenue*}
  \erdbang{brands}{2,0}{brand_id:PK, brand_name}
  \erdnoi[kieu=n-1, tren={\code{campaigns.brand_id = brands.brand_id}}]{campaigns.brand_id}{brands.brand_id}
\end{erd}
\begin{nghiepvu}{Chọn chiến dịch dịp 8/3 tiêu biểu của từng thương hiệu để chia sẻ cách làm trong nội bộ.}
  \khainiem{chien-dich}
  \luuy{``Cao nhất của mỗi thương hiệu'' nghĩa là so trong nội bộ từng thương hiệu; mức cao nhất của cả ba thương hiệu không dùng ở đây. Phạm vi ngày phải giống nhau ở truy vấn ngoài và truy vấn con.}
\end{nghiepvu}
\goiy{Truy vấn con tính \code{MAX(revenue)} cần biết đang xét thương hiệu nào. Lấy thông tin đó từ dòng nào?}
\goiy{\code{WHERE c.revenue = (SELECT MAX(c2.revenue) FROM campaigns c2 WHERE c2.brand_id = c.brand_id AND c2.start_date = '2025-03-05')}.}
\ghichuthietke{Dễ: 2 bảng, chữ trên đường nối. Phạm vi một ngày 5/3/2025 (bản đồ ghi cả kỳ) để truy vấn con tương quan chạy nhanh. Sửa lần 1 sau kiểm B: bản 8/3 ra cùng kết quả ở hai biến thể và bối cảnh ``rất nhiều chiến dịch'' sai số liệu. Sửa lần 2: bản tuần 1--8/3 chạy hơn 200 giây trên hidden-1 (truy vấn con dò lại khoảng ngày cho từng dòng), nên thu về một ngày; đo lúc máy tải: lời giải 0,8 giây (public) và 1,3 giây (hidden-1), cách viết khoảng nửa mở trong truy vấn con khoảng 2 giây. Ngày 5/3 có NY-CMP-53612 nên hidden-1 khác public ở doanh thu (biến thể c). Bẫy ở public: truy vấn con không tương quan ra 1 dòng; gom nhóm lấy \code{MIN(campaign_id)} kèm \code{MAX(revenue)} ra mã không khớp doanh thu. Public không có hoà ở mức cao nhất.}
\end{cauhoi}

% ================================================================ Q3 · Viết · Cơ bản · C3 · marketing
\begin{cauhoi}[ma=mkt-q3-social-media-purplle-nhap-tren-binh-quan, dulieu=mkt, domain=marketing,
  hinhthuc=write, cap=basic, dang=C3, kynang={truy vấn con một giá trị}, bay={binh-quan-sai-pham-vi},
  tacgia={A}, trangthai=review, bokiemtra={public, hidden-1}, thutu=co, socot=vitri]
\tieude{Chiến dịch Social Media nào của Purplle có lượt nhấp trên bình quân?}
\boicanh[Chị Ngân · Quản lý thương hiệu Purplle]{Ngày 6/5 Purplle mở khá nhiều chiến dịch marketing mạng xã hội. Em lọc giúp chị những chiến dịch có lượt nhấp cao hơn bình quân lượt nhấp của mọi chiến dịch trong dữ liệu.}
\begin{yeucau}
  \item Xét các chiến dịch của Purplle (\code{brand_id = 2}), loại Social Media (marketing mạng xã hội, \code{type_id = 5}), bắt đầu ngày 6/5/2025.
  \item Lọc các chiến dịch có lượt nhấp (click, cột \code{clicks}) lớn hơn bình quân \code{clicks} của mọi chiến dịch trong bảng \code{campaigns} (cả ba thương hiệu, cả kỳ dữ liệu).
  \item Trả về \code{campaign_id}, \code{clicks}; sắp \code{clicks} giảm dần, bằng nhau thì \code{campaign_id} tăng dần.
\end{yeucau}
\begin{erd}
  \erdbang{campaigns}{0,0}{campaign_id:PK='PU-CMP-10990', brand_id*=2, type_id*=5, start_date*=2025-05-05, clicks*=7835}
\end{erd}
\begin{nghiepvu}{Chọn các chiến dịch mạng xã hội thu hút nhiều lượt nhấp để xem lại nội dung quảng cáo.}
  \khainiem{chien-dich}
  \luuy{Bình quân dùng để so là một con số chung cho mọi dòng, tính trên toàn bảng chứ không chỉ trên các chiến dịch đang xét.}
\end{nghiepvu}
\goiy{\code{WHERE clicks > (SELECT AVG(clicks) FROM campaigns)}, cùng các điều kiện thương hiệu, loại và ngày.}
\ghichuthietke{Cơ bản (v4, hạ từ Trung bình). Năm trục: bảng: 1; bước: 1 tầng (lọc với một truy vấn con một giá trị, truy vấn con không lọc); mơ hồ: yêu cầu ghi đủ cột, giá trị lọc và phạm vi bình quân; bẫy: 0; nghiệp vụ: một khái niệm. Đúng mẫu \code{q1-tren-gia-trung-binh} của lộ trình. Đạt Cơ bản 5/5. ERD 1 bảng, giá trị mẫu lấy từ ngày 5/5 (không thuộc kết quả), không bảng nhiễu. Bản Trung bình cũ (CTR của từng loại so với CTR chung, HAVING) đã bỏ. Kết quả public 10 dòng, hidden-1 8 dòng. Wrong: \code{binh_quan_cung_nhom} (bình quân của chính nhóm Purplle Social Media ngày 6/5, public ra 13 dòng), \code{thieu_loc_loai} (lấy cả năm loại); cả hai rớt public. Ô: C3 · marketing · basic · write, không trùng trong gói.}
\end{cauhoi}

% ================================================================ Q4 · Viết · Khó · C3 · finance
\begin{cauhoi}[ma=mkt-q4-thang-roi-tren-muc-quy-1, dulieu=mkt, domain=finance,
  hinhthuc=write, cap=hard, dang=C3, kynang={truy vấn con tương quan, bảng dẫn xuất, tỷ lệ của tổng}, bay={avg-cua-ty-le, khong-tuong-quan, muc-chung-sai-pham-vi, lam-tron-truoc-khi-so},
  tacgia={A}, trangthai=review, bokiemtra={public, hidden-1}, thutu=co, socot=vitri, saiso=0.0001]
\tieude{Tháng nào quý 1/2025 thương hiệu đạt ROI trên mức cả quý?}
\boicanh[Anh Đức · Giám đốc tài chính]{Anh đang soạn báo cáo hiệu quả chi phí marketing quý 1/2025 cho ban giám đốc. Ở mỗi thương hiệu, anh muốn biết tháng nào có ROI tốt hơn mức của chính thương hiệu đó trong cả quý.}
\begin{yeucau}
  \item Xét các chiến dịch bắt đầu từ 1/1/2025 tới hết 31/3/2025.
  \item Với mỗi thương hiệu, lấy các tháng có ROI (lợi tức đầu tư) của tháng cao hơn ROI của chính thương hiệu đó trong cả quý 1/2025. ROI của tháng và của quý tính cho cả nhóm chiến dịch theo khái niệm ROI ở tab Nghiệp vụ; so trên giá trị chưa làm tròn.
  \item Trả về \code{brand_id}, \code{thang} (chuỗi dạng \code{'2025-01'}), \code{roi_thang}, \code{roi_quy} (hai cột sau làm tròn 4 chữ số); sắp \code{brand_id} rồi \code{thang} tăng dần.
\end{yeucau}
\begin{erd}
  \erdbang{campaigns}{1,0}{campaign_id:PK, brand_id:FK, type_id:FK, start_date:FK}
  \erdbang{brands}{0,0}{brand_id:PK}
  \erdbang{campaign_types}{2,0}{type_id:PK}
  \erdbang{calendar}{0,1}{cal_date:PK}
  \erdbang{campaign_channels}{2,1}{campaign_id:FK, channel_id:FK}
  \erdkhoa{Lời giải chỉ cần \code{campaigns}: tháng lấy từ \code{start_date}, thương hiệu từ \code{brand_id}. \code{brands}, \code{campaign_types}, \code{calendar}, \code{campaign_channels} không cần.}
\end{erd}
\begin{nghiepvu}{Chọn tháng làm mốc khi rà soát cách chi tiêu marketing của từng thương hiệu cho quý sau.}
  \khainiem{roi}
  \khainiem{chi-phi-chien-dich}
  \khainiem{binh-quan-gia-quyen}
  \luuy{Mức so sánh của mỗi thương hiệu là ROI của chính thương hiệu đó trong quý 1/2025, không phải ROI chung của ba thương hiệu hay ROI của cả kỳ dữ liệu. Đây là dữ liệu tổng hợp nên ROI các tháng chênh nhau rất ít; vì vậy phải so trên giá trị chưa làm tròn.}
\end{nghiepvu}
\goiy{Mỗi dòng cần hai con số: ROI của tháng (gom theo thương hiệu và tháng) và ROI của cả quý (gom theo thương hiệu). Truy vấn con tính ROI quý phải gắn với thương hiệu của dòng đang xét.}
\goiy{ROI của một nhóm là \code{(SUM(revenue) - SUM(acquisition_cost * conversions)) / SUM(acquisition_cost * conversions)}.}
\ghichuthietke{Khó (v4, viết mới thay câu Trung bình \code{mkt-q4-doi-soat-roi-tinh-san}). Năm trục: bảng: lời giải 1 bảng (mẫu o4 Khó dùng 1 bảng); ERD 5 bảng chỉ khoá, 4 bảng nhiễu cho đủ 5 bảng, cách nối ở \code{erdkhoa}; bước: 4 tầng (ROI theo thương hiệu và tháng, ROI theo thương hiệu cả quý bằng truy vấn con tương quan hoặc bảng dẫn xuất, so, làm tròn); mơ hồ: lời nhắn nói bằng ngôn ngữ báo cáo, công thức ROI nhóm nằm ở tab Nghiệp vụ; bẫy: 4 (bình quân cột \code{roi}; mức chung của ba thương hiệu; mức của cả kỳ; làm tròn trước khi so, chỉ lộ ở hidden-1: tháng 1 của Tira Beauty cao hơn ROI quý chưa tới 0,001); nghiệp vụ: ROI của tổng với chi phí tự tính, so trong từng thương hiệu. Đạt Khó 4/5. Thu phạm vi về một quý: truy vấn con tương quan đặt ở cả SELECT và HAVING chạy khoảng 4 giây lúc máy tải (bản 5 tháng đầu năm mất 7,5 giây). \code{saiso=0.0001} vì ROI tháng chênh ROI quý dưới 0,01. Kết quả public 6 dòng, hidden-1 6 dòng khác giá trị. Wrong: \code{binh_quan_cot_roi}, \code{muc_chung_ca_ba_thuong_hieu}, \code{muc_so_sanh_ca_ky} (rớt public), \code{lam_tron_truoc_khi_so} (chỉ rớt hidden-1). Ô: C3 · finance · hard · write, không trùng trong gói (q1 cùng dạng, cùng mảng nhưng cấp Dễ).}
\end{cauhoi}

% ================================================================ Q5 · Viết · Dễ · C2 · sales
\begin{cauhoi}[ma=mkt-q5-ty-trong-loai-purplle, dulieu=mkt, domain=sales,
  hinhthuc=write, cap=easy, dang=C2, kynang={truy vấn con một giá trị, tỷ trọng}, bay={sai-mau-so, thieu-loc-ngoai},
  tacgia={A}, trangthai=review, bokiemtra={public, hidden-1}, thutu=co, socot=vitri, saiso=0.01]
\tieude{Mỗi loại chiến dịch đóng góp bao nhiêu phần trăm doanh thu Purplle?}
\boicanh[Chị Ngân · Quản lý thương hiệu Purplle]{Chị đang viết báo cáo cơ cấu doanh thu của Purplle cả kỳ. Chị cần doanh thu và tỷ trọng doanh thu của từng loại chiến dịch trong tổng doanh thu Purplle.}
\begin{yeucau}
  \item Xét mọi chiến dịch của Purplle (\code{brand_id = 2}) trong kỳ dữ liệu.
  \item Tính tổng doanh thu và tỷ trọng của từng loại chiến dịch, với tỷ trọng = doanh thu của loại × 100 / tổng doanh thu của Purplle.
  \item Trả về \code{type_name}, \code{doanh_thu}, \code{ty_trong_pct} (phần trăm, làm tròn 2 chữ số); sắp \code{doanh_thu} giảm dần.
\end{yeucau}
\begin{erd}
  \erdbang{campaigns}{0,0}{campaign_id:PK, brand_id*, type_id:FK, revenue*}
  \erdbang{campaign_types}{2,0}{type_id:PK, type_name*}
  \erdnoi[kieu=n-1, tren={\code{campaigns.type_id = campaign_types.type_id}}]{campaigns.type_id}{campaign_types.type_id}
\end{erd}
\begin{nghiepvu}{Mô tả cơ cấu doanh thu của Purplle theo loại chiến dịch trong báo cáo thương hiệu.}
  \khainiem{chien-dich}
  \luuy{Tỷ trọng của các phần trong một tổng cộng lại bằng 100\%. Mẫu số phải là tổng của đúng phạm vi đang chia nhỏ, ở đây là riêng Purplle.}
\end{nghiepvu}
\goiy{Tổng doanh thu của Purplle là một con số dùng chung cho mọi dòng; có thể lấy bằng một truy vấn con.}
\goiy{\code{SUM(c.revenue) * 100 / (SELECT SUM(revenue) FROM campaigns WHERE brand_id = 2)}.}
\ghichuthietke{Dễ (hạ từ Trung bình sau kiểm C: một truy vấn con vô hướng cộng GROUP BY, cùng dạng q1): 2 bảng, chữ trên đường nối. Bẫy ở public: mẫu số là doanh thu cả ba thương hiệu (các tỷ trọng cộng lại khoảng một phần ba); quên lọc Purplle ở truy vấn ngoài (tử số lấy cả ba thương hiệu, tổng khoảng 300 phần trăm). Bỏ bài sai chia cho bình quân một loại (lỗi ít gặp). Tỷ trọng các loại gần bằng nhau (dữ liệu tổng hợp), nên sắp theo doanh thu để tránh hoà khi làm tròn. hidden-1 bỏ chiến dịch đuôi 7.}
\end{cauhoi}

% ================================================================ Q6 · Viết · Trung bình · A1 · marketing
\begin{cauhoi}[ma=mkt-q6-ngay-purplle-khong-seo, dulieu=mkt, domain=marketing,
  hinhthuc=write, cap=medium, dang=A1, kynang={NOT EXISTS, truy vấn con tương quan}, bay={khong-tuong-quan, thieu-can-tren, nham-doi-tuong-voi-phan-khuc},
  tacgia={A}, trangthai=review, bokiemtra={public, hidden-1}, thutu=co, socot=vitri]
\tieude{Ngày nào tháng 5/2025 Purplle không mở chiến dịch SEO cho giới trẻ?}
\boicanh[Chị Hằng · Chuyên viên SEO]{Chị đang rà nhịp mở chiến dịch SEO (tối ưu hoá công cụ tìm kiếm) tiếng Tamil nhắm giới trẻ (Youth) của Purplle. Tháng 5 có những ngày nào không có chiến dịch nào như vậy khởi chạy?}
\begin{yeucau}
  \item Xét mọi ngày từ 1/5/2025 tới hết 31/5/2025 trong lịch ngày.
  \item Lấy các ngày không có chiến dịch nào của Purplle (\code{brand_id = 2}), loại SEO (\code{type_id = 4}), ngôn ngữ Tamil (\code{language_id = 4}), đối tượng mục tiêu Youth (\code{target_audience_id = 5}) bắt đầu vào ngày đó.
  \item Trả về \code{cal_date}, \code{day_name}; sắp \code{cal_date} tăng dần.
\end{yeucau}
\begin{erd}
  \erdbang{calendar}{0,0}{cal_date:PK, day_name*, is_weekend}
  \erdbang{campaigns}{2,0}{campaign_id:PK, start_date:FK*, brand_id*, type_id*, language_id*, target_audience_id*, segment_id}
  \erdbang{campaign_types}{2,1}{type_id:PK, type_name}
  \erdnoi[kieu=1-n]{calendar.cal_date}{campaigns.start_date}
  \erdnoi[kieu=n-1]{campaigns.type_id}{campaign_types.type_id}
\end{erd}
\begin{nghiepvu}{Tìm những ngày SEO tiếng Tamil cho giới trẻ của Purplle bị bỏ trống để lên lịch bù.}
  \khainiem{doi-tuong-va-phan-khuc}
  \luuy{Ngày cần xét lấy từ lịch ngày. Ngày không có chiến dịch phù hợp thì không có dòng nào để lọc trong \code{campaigns}; phải hỏi ngược ``có tồn tại chiến dịch nào vào ngày này không''. Tháng 5/2025 là tháng đủ dữ liệu.}
\end{nghiepvu}
\goiy{Điều kiện ``không tồn tại chiến dịch thoả điều kiện vào ngày này'' viết bằng \code{NOT EXISTS}; truy vấn con phải gắn với ngày đang xét của lịch.}
\goiy{\code{WHERE NOT EXISTS (SELECT 1 FROM campaigns c WHERE c.brand_id = 2 AND c.type_id = 4 AND c.language_id = 4 AND c.target_audience_id = 5 AND c.start_date = k.cal_date)}.}
\ghichuthietke{Trung bình: 3 bảng, ký hiệu 1-n. Sửa sau kiểm B: bản cũ (Purplle SEO tháng 6) ra 7 ngày cuối tháng sau khi hết dữ liệu, giống nhau ở hai biến thể. Bản mới chọn tháng đủ và thu hẹp nhóm: public 10 ngày, hidden-1 thêm ngày 30/5. Bẫy ở public: truy vấn con không tương quan ra 0 dòng; lịch thiếu cận trên lấy cả tháng 6; lọc Youth theo \code{segment_id}. Bỏ bài sai lấy ngày từ \code{campaigns} vì tháng 5 ngày nào cũng có chiến dịch nên ra trùng lời giải. NOT IN cho cùng kết quả vì \code{start_date} không NULL.}
\end{cauhoi}

% ================================================================ Q7 · Sửa lỗi · Cơ bản · A1 · customer
\begin{cauhoi}[ma=mkt-q7-shoppers-ra-rong, dulieu=mkt, domain=customer,
  hinhthuc=debug, cap=basic, dang=A1, kynang={IN (SELECT ...), LIKE}, bay={mau-like-sai-vi-tri},
  tacgia={A}, trangthai=review, bokiemtra={public, hidden-1}, thutu=co, socot=vitri]
\tieude{Chiến dịch nhắm nhóm ``Shoppers'' không ra dòng nào}
\boicanh[Bạn Linh · Thực tập sinh]{Em đếm chiến dịch tháng 5/2025 nhắm nhóm khách mua sắm (tên nhóm có chữ ``Shoppers'') theo từng thương hiệu mà bảng không ra dòng nào, dù tháng nào phòng mình cũng chạy chiến dịch cho nhóm khách này.}
\begin{yeucau}
  \item Mục đích của truy vấn khởi đầu: đếm số chiến dịch của từng thương hiệu bắt đầu từ 1/5/2025 tới hết 31/5/2025 mà đối tượng mục tiêu có tên chứa chữ \code{Shoppers}.
  \item Sửa truy vấn cho đúng yêu cầu; giữ cột \code{brand_id}, \code{so_chien_dich}, sắp \code{brand_id} tăng dần.
\end{yeucau}
\begin{erd}
  \erdbang{campaigns}{0,0}{campaign_id:PK, target_audience_id:FK*, start_date*}
  \erdbang{audiences}{2,0}{audience_id:PK, audience_name*='Premium Shoppers'}
  \erdnoi[kieu=n-1, tren={đối tượng mục tiêu}]{campaigns.target_audience_id}{audiences.audience_id}
\end{erd}
\begin{nghiepvu}{Đếm chiến dịch nhắm khách mua sắm cao cấp của từng thương hiệu để so mức đầu tư vào nhóm này.}
  \khainiem{doi-tuong-va-phan-khuc}
  \luuy{Trong mẫu \code{LIKE}, dấu \code{\%} thay cho một chuỗi bất kỳ. Vị trí của \code{\%} quyết định chữ cần tìm nằm ở đầu, ở cuối hay ở giữa tên.}
\end{nghiepvu}
\goiy{Chạy riêng truy vấn con: ra mấy mã nhóm? Tên nhóm có chữ Shoppers bắt đầu bằng chữ gì?}
\ghichuthietke{Cơ bản (hạ từ Dễ sau kiểm C: lỗi một chỗ ở giá trị lọc, giá trị mẫu là manh mối trực tiếp). 2 bảng vì lỗi nằm trong truy vấn con trên \code{audiences}; giá trị mẫu \code{'Premium Shoppers'} là manh mối, không tô cảnh báo đường nối (đường nối đúng). Bẫy ở public: sửa mẫu nhưng viết lại điều kiện ngày thiếu cận trên; sửa mẫu nhưng lọc theo \code{segment_id}. Bỏ bài sai \code{= 'Shoppers'} vì ra trùng truy vấn khởi đầu. hidden-1 bỏ chiến dịch đuôi 7.}
\end{cauhoi}

% ================================================================ Q8 · Sửa lỗi · Trung bình · E2 · operations
\begin{cauhoi}[ma=mkt-q8-dong-sach-khong-nhieu-sai, dulieu=mkt, domain=operations,
  hinhthuc=debug, cap=medium, dang=E2, kynang={NOT EXISTS, khoá kép}, bay={khoa-kep, row-index-lap-giua-thuong-hieu},
  tacgia={A}, trangthai=review, bokiemtra={public, hidden-1}, thutu=co, socot=vitri]
\tieude{Số dòng dữ liệu thô không có trong nhật ký nhiễu bị sai}
\boicanh[Anh Khánh · Chuyên viên phân tích marketing]{Anh cần biết mỗi file thô có bao nhiêu dòng không bị ghi trong nhật ký nhiễu. Bảng hiện tại ra số thấp hơn anh ước lượng khá nhiều.}
\begin{yeucau}
  \item Truy vấn khởi đầu phải đếm số dòng \code{campaigns_raw} của từng thương hiệu không bị nhật ký nhiễu nhắc tới, tức không có dòng \code{noise_log} nào cùng thương hiệu và cùng vị trí dòng.
  \item Sửa truy vấn cho đúng yêu cầu; giữ cột \code{brand_id}, \code{so_dong}, sắp \code{brand_id} tăng dần.
\end{yeucau}
\begin{erd}
  \erdbang{campaigns_raw}{0,0}{raw_id:PK, brand_id:FK, row_index, campaign_id}
  \erdbang{noise_log}{2,0}{log_id:PK, brand_id:FK, row_index, campaign_id, noise_type}
  \erdbang{brands}{1,1}{brand_id:PK, brand_name}
  \erdnoi[kieu=n-1]{campaigns_raw.brand_id}{brands.brand_id}
  \erdnoi[kieu=n-1]{noise_log.brand_id}{brands.brand_id}
\end{erd}
\begin{nghiepvu}{Biết bao nhiêu dòng của dữ liệu thô không bị nhật ký nhiễu nhắc tới, để ước lượng khối lượng phải kiểm tra.}
  \quytrinh{Làm sạch dữ liệu}{nạp, kiểm tra, sửa, đối soát}{}
  \khainiem{du-lieu-tho}
  \khainiem{nhat-ky-nhieu}
  \luuy{\code{row_index} đếm lại từ 0 trong mỗi file, nên cùng một vị trí dòng có ở cả ba thương hiệu. Một dòng của dữ liệu thô chỉ được xác định khi biết cả thương hiệu và vị trí. Với dòng trùng, nhật ký ghi vị trí của dòng gốc, nên bản chép không bị nhắc tới; dòng không có trong nhật ký chưa chắc đã sạch.}
\end{nghiepvu}
\goiy{Vị trí dòng 120 bị ghi nhiễu ở file Nykaa thì dòng 120 của file Purplle có bị loại theo truy vấn hiện tại không?}
\goiy{Dùng \code{NOT EXISTS} với hai điều kiện \code{n.brand_id = r.brand_id AND n.row_index = r.row_index}.}
\ghichuthietke{Trung bình: 3 bảng, ký hiệu n-1, không vẽ đường nối giữa hai bảng nghiệp vụ (khoá dòng kép nằm trong khái niệm). Bẫy ở public: so theo \code{campaign_id} loại luôn bản chép của dòng trùng; NOT EXISTS theo \code{brand_id} và \code{campaign_id} (cùng lỗi theo cách khác). Sửa sau kiểm B: bỏ bài sai chỉ so \code{row_index} vì ra trùng truy vấn khởi đầu; tab Nghiệp vụ không còn gọi các dòng này là dùng được ngay (trong đó có 50 bản chép trùng mỗi file). hidden-1 bỏ dòng dữ liệu thô và nhật ký của chiến dịch đuôi 7, thêm một dòng dữ liệu thô.}
\end{cauhoi}

% ================================================================ Q9 · Sửa lỗi · Khó · C3 · marketing
\begin{cauhoi}[ma=mkt-q9-tuong-tac-tren-binh-quan-hai-loi, dulieu=mkt, domain=marketing,
  hinhthuc=debug, cap=hard, dang=C3, kynang={truy vấn con tương quan, bí danh}, bay={tuong-quan-sai-bi-danh, binh-quan-sai-pham-vi},
  tacgia={A}, trangthai=review, bokiemtra={public, hidden-1}, thutu=co, socot=vitri]
\tieude{Chiến dịch tương tác tốt dịp 8/3 của Nykaa bị đếm sai}
\boicanh[Chị Trang · Quản lý thương hiệu Nykaa]{Chị muốn biết trong hai ngày 7 và 8/3, mỗi loại chiến dịch của Nykaa có bao nhiêu chiến dịch có điểm tương tác cao hơn bình quân của chính loại đó. Bạn thực tập đã đếm thử, nhưng chị thấy các con số chưa khớp với bình quân từng loại. Em kiểm tra lại giúp chị.}
\begin{yeucau}
  \item Cần đếm số chiến dịch của Nykaa (\code{brand_id = 1}) theo từng loại, xét các chiến dịch bắt đầu từ 7/3/2025 tới hết 8/3/2025 có \code{engagement_score} lớn hơn bình quân \code{engagement_score} của các chiến dịch Nykaa \textbf{cùng loại} trong cùng hai ngày đó.
  \item Truy vấn khởi đầu có hơn một chỗ sai; sửa một chỗ vẫn chưa đúng.
  \item Giữ cột \code{type_id}, \code{so_chien_dich}, sắp \code{type_id} tăng dần.
\end{yeucau}
\begin{erd}
  \erdbang{campaigns}{1,0}{campaign_id:PK, brand_id:FK, type_id:FK, target_audience_id:FK, start_date:FK}
  \erdbang{campaign_types}{2,0}{type_id:PK}
  \erdbang{brands}{0,0}{brand_id:PK}
  \erdbang{audiences}{0,1}{audience_id:PK}
  \erdbang{calendar}{2,1}{cal_date:PK}
  \erdkhoa{Lời giải chỉ cần \code{campaigns}; truy vấn con đọc lại chính bảng này. \code{campaign_types}, \code{brands}, \code{audiences}, \code{calendar} không cần.}
\end{erd}
\begin{nghiepvu}{Tìm chiến dịch có nội dung thu hút hơn mức thường thấy của cùng loại trong dịp 8/3 để làm mẫu cho đợt sau.}
  \khainiem{diem-tuong-tac}
  \luuy{Mỗi chiến dịch được so với bình quân của các chiến dịch cùng thương hiệu, cùng loại và cùng hai ngày; không so với bình quân của mọi loại hay của cả ba thương hiệu.}
\end{nghiepvu}
\goiy{Trong truy vấn con, tên cột không ghi bí danh được hiểu là cột của bảng gần nhất. Vậy \code{type_id = type_id} so cột nào với cột nào? Truy vấn con đang lấy chiến dịch của những thương hiệu nào?}
\goiy{Đặt bí danh \code{campaigns c2} trong truy vấn con và lọc \code{c2.brand_id = 1 AND c2.type_id = c.type_id} cùng hai ngày.}
\ghichuthietke{Khó (v4, nâng từ Trung bình). Năm trục: bảng: lời giải 1 bảng đọc hai lần; ERD 5 bảng chỉ khoá, 4 bảng nhiễu cho đủ 5 bảng (mẫu d4, o4), không vẽ đường nối, cách nối ở \code{erdkhoa}; sửa lỗi: hai lỗi chồng nhau (\code{type_id = type_id} không tương quan; truy vấn con thiếu lọc Nykaa), sửa một lỗi vẫn sai; bẫy: 3 (thiếu tương quan, bình quân lấy cả ba thương hiệu, bỏ lọc Nykaa ở truy vấn ngoài khi sửa); mơ hồ: yêu cầu không chỉ chỗ sai, lời nhắn chỉ nêu triệu chứng, ERD trung tính; nghiệp vụ: so trong nhóm theo ba điều kiện (thương hiệu, loại, kỳ). Đạt Khó 3/5 cộng trục sửa lỗi. Thu phạm vi từ 7 ngày (1--7/3) còn 2 ngày (7--8/3): bản 7 ngày lời giải mất 7,7 giây và bài sai chỉ sửa bí danh mất 23,5 giây lúc máy tải; bản 2 ngày lời giải dưới 1 giây, bài sai dưới 2,1 giây. Kết quả public khác hidden-1 (bỏ chiến dịch đuôi 7). Wrong: \code{chi_sua_bi_danh} (bình quân loại của cả ba thương hiệu), \code{chi_them_loc_thuong_hieu} (bình quân mọi loại của Nykaa); cả hai rớt public. Ô: C3 · marketing · hard · debug, không trùng trong gói (q3 cùng dạng, cùng mảng nhưng cấp Cơ bản, hình thức Viết).}
\end{cauhoi}


% ################################################################ CASE STUDY

% ================================================================ C5 · Case · Dễ · C3 · marketing
\begin{cauhoi}[ma=mkt-c5-kenh-hieu-qua-hon, dulieu=mkt, domain=marketing,
  hinhthuc=case, cap=easy, diem=20, tacgia={A}, trangthai=review, bokiemtra={public, hidden-1}, thutu=co, saiso=0.01]
\tieude{Kênh nào gắn với chiến dịch hiệu quả hơn?}
\boicanh[Anh Vinh · Trưởng nhóm quảng cáo trả phí]{Phòng mình đang tính dồn thêm ngân sách nửa cuối năm cho một vài kênh. Anh cần biết trong nửa đầu 2025, chiến dịch có dùng kênh nào có tỷ lệ nhấp cao hơn và ROI tốt hơn, kèm một kết luận ngắn.}
\begin{erd}
  \erdbang[hl]{campaigns}{0,0}{campaign_id:PK, start_date*, impressions*, clicks*, revenue*, acquisition_cost*, conversions*}
  \erdbang[hl]{campaign_channels}{1,0}{campaign_id:FK, channel_id:FK}
  \erdbang[hl]{channels}{2,0}{channel_id:PK, channel_name*}
  \erdnoi[kieu=1-n, tren={một chiến dịch, 1 tới 3 kênh}]{campaigns.campaign_id}{campaign_channels.campaign_id}
  \erdnoi[kieu=n-1, tren={nhiều dòng, một kênh}]{campaign_channels.channel_id}{channels.channel_id}
\end{erd}
\begin{nghiepvu}{Có nên dồn ngân sách nửa cuối năm cho một kênh cụ thể hay không.}
  \khainiem{kenh-da-tri}
  \khainiem{ctr}
  \khainiem{roi}
  \luuy{Một chiến dịch chạy nhiều kênh nên doanh thu không chia được cho từng kênh. Số liệu ``theo kênh'' là số liệu của các chiến dịch có dùng kênh đó, và các nhóm này chồng lên nhau.}
\end{nghiepvu}
\begin{buoc}[ma=b1, diem=8, hinhthuc=write, bang={campaigns, campaign_channels, channels}]{Số chiến dịch và CTR theo kênh}
  \begin{yeucau}
    \item Xét các chiến dịch của ba thương hiệu bắt đầu từ 1/1/2025 tới hết 30/6/2025. Chiến dịch chạy nhiều kênh được tính ở mọi kênh mà chiến dịch đó dùng.
    \item Đếm số chiến dịch có dùng mỗi kênh và tính CTR (tỷ lệ nhấp) của kênh đó theo tỷ lệ của tổng.
    \item Trả về \code{channel_name}, \code{so_chien_dich}, \code{ctr_pct} (phần trăm, làm tròn 2 chữ số); sắp \code{ctr_pct} (giá trị đã làm tròn) giảm dần, bằng nhau thì \code{channel_name} tăng dần.
  \end{yeucau}
  \goiy{CTR của một kênh là \code{ROUND(SUM(c.clicks) * 100 / SUM(c.impressions), 2)} trên các chiến dịch dùng kênh đó.}
\end{buoc}
\begin{buoc}[ma=b2, diem=7, hinhthuc=write, bang={campaigns, campaign_channels, channels}]{ROI tổng theo kênh}
  \begin{yeucau}
    \item Cùng phạm vi như bước 1. Ở mỗi kênh, tính tổng doanh thu, tổng chi phí (chi phí mỗi chiến dịch = \code{acquisition_cost} × \code{conversions}) và ROI (lợi tức đầu tư) theo tỷ lệ của tổng.
    \item Trả về \code{channel_name}, \code{doanh_thu}, \code{chi_phi} (2 chữ số), \code{roi} (4 chữ số); sắp \code{roi} (giá trị đã làm tròn) giảm dần, bằng nhau thì \code{channel_name} tăng dần.
  \end{yeucau}
  \goiy{ROI của nhóm = (tổng doanh thu - tổng chi phí) / tổng chi phí; không lấy \code{AVG(roi)}.}
\end{buoc}
\begin{buoc}[ma=b3, diem=5, hinhthuc=ketluan, bang={campaigns, channels}]{Kết luận}
  \begin{luachon}
    \phuongan*{Các kênh chênh nhau rất ít về CTR và ROI, và một chiến dịch được tính ở nhiều kênh, nên chưa đủ căn cứ để dồn ngân sách cho một kênh.}
      {Đúng: Bước 1 và 2 cho thấy các kênh chênh nhau rất nhỏ, và một chiến dịch được tính ở nhiều kênh nên không tách được tác động từng kênh.}
    \phuongan{Dù CTR các kênh gần nhau, cộng doanh thu sáu kênh ở bước 2 ra tổng doanh thu nửa đầu 2025, nên kênh có doanh thu lớn nhất đáng được dồn ngân sách.}
      {Chiến dịch nhiều kênh bị cộng ở mỗi kênh, nên tổng sáu kênh lớn hơn doanh thu thật; doanh thu không quy được cho riêng một kênh.}
    \phuongan{Email nhỉnh hơn các kênh khác ở cả CTR lẫn ROI dù chênh lệch nhỏ, nên có thể chuyển một phần ngân sách từ Facebook sang Email.}
      {ROI theo kênh chỉ là ROI của nhóm chiến dịch có dùng kênh đó, chưa đo tác động riêng của kênh; chênh lệch lại quá nhỏ để kết luận.}
  \end{luachon}
\end{buoc}
\ghichuthietke{Case Dễ: 3 bảng có chữ trên đường nối, mọi bước dùng cùng bộ bảng. Trên public CTR theo kênh khoảng 8,46--8,53\%, ROI khoảng 1,92--1,98; Email đứng đầu cả hai nhưng chênh rất nhỏ; Instagram và WhatsApp hoà CTR khi làm tròn (quy tắc hoà theo tên kênh). Bài sai bước 1: bình quân CTR từng chiến dịch (public), thiếu cận trên (hidden-1: chiến dịch 2025-07-01 cộng vào WhatsApp, YouTube). Bước 2: \code{AVG(roi)}, coi \code{acquisition_cost} là chi phí cả chiến dịch. Kết luận không khẳng định nguyên nhân (STYLE-mkt B.6, M4, M5). Sửa sau kiểm B: đề ghi rõ sắp theo giá trị đã làm tròn (sắp theo giá trị chưa làm tròn đổi chỗ Instagram và WhatsApp); cân độ dài phương án, phương án sai cũng có từ rào (K32).}
\end{cauhoi}

% ================================================================ C6 · Case · Trung bình · E2 · operations
\begin{cauhoi}[ma=mkt-c6-tho-va-sach-lech-o-dau, dulieu=mkt, domain=operations,
  hinhthuc=case, cap=medium, diem=30, tacgia={A}, trangthai=review, bokiemtra={public, hidden-1}, thutu=co]
\tieude{Bản thô và bản sạch lệch nhau ở đâu?}
\boicanh[Chị Hạnh · Trưởng phòng marketing số]{Đội dữ liệu gửi cả bản thô lẫn bản sạch của ba file chiến dịch. Trước khi cho các nhóm dùng bản sạch, chị cần biết hai bản lệch nhau ở đâu, lệch bao nhiêu, và nhật ký nhiễu có giải thích được hết không.}
\begin{erd}
  \erdbang[hl]{campaigns_raw}{0,0}{raw_id:PK, brand_id:FK, row_index, campaign_id, impressions, clicks, leads, conversions, revenue, acquisition_cost}
  \erdbang[hl]{campaigns}{2,0}{campaign_id:PK, brand_id:FK, impressions, clicks, leads, conversions, revenue, acquisition_cost}
  \erdbang[hl]{noise_log}{0,1}{log_id:PK, brand_id:FK, row_index, campaign_id, noise_type}
  \erdnoi[kieu=n-1]{campaigns_raw.campaign_id}{campaigns.campaign_id}
  \erdnoi[kieu=1-n]{campaigns_raw.row_index}{noise_log.row_index}
\end{erd}
\begin{nghiepvu}{Cho phép các nhóm dùng bản sạch hay yêu cầu đội dữ liệu kiểm tra lại quy trình làm sạch.}
  \quytrinh{Làm sạch dữ liệu}{nạp, kiểm tra, sửa, đối soát}{}
  \khainiem{du-lieu-tho}
  \khainiem{nhat-ky-nhieu}
  \khainiem{dong-trung}
  \luuy{Khoá dòng của bản thô là cặp (\code{brand_id}, \code{row_index}); \code{campaign_id} lặp lại ở dòng trùng. So một ô có thể trống với một ô có giá trị cần phép so an toàn với NULL (\code{<=>}), vì \code{NULL <> 5} cho NULL chứ không cho đúng, nên dòng đó bị bỏ qua.}
\end{nghiepvu}
\begin{buoc}[ma=b1, diem=7, hinhthuc=write, bang={campaigns_raw, campaigns}]{Số dòng và số mã của hai bản}
  \begin{yeucau}
    \item Theo từng thương hiệu, đếm số dòng của bản thô, số \code{campaign_id} khác nhau của bản thô và số chiến dịch của bản sạch.
    \item Trả về \code{brand_id}, \code{so_dong_tho}, \code{so_ma_tho}, \code{so_chien_dich_sach}; sắp \code{brand_id} tăng dần.
  \end{yeucau}
  \goiy{Số chiến dịch của bản sạch lấy bằng truy vấn con tương quan \code{(SELECT COUNT(*) FROM campaigns c WHERE c.brand_id = r.brand_id)}.}
\end{buoc}
\begin{buoc}[ma=b2, diem=8, hinhthuc=write, bang={campaigns_raw, campaigns}]{Dòng thô có chỉ số khác bản sạch}
  \begin{yeucau}
    \item Nối mỗi dòng của bản thô với chiến dịch cùng \code{campaign_id} ở bản sạch.
    \item Đếm, theo thương hiệu, số dòng của bản thô có ít nhất một trong sáu cột \code{impressions}, \code{clicks}, \code{leads}, \code{conversions}, \code{revenue}, \code{acquisition_cost} khác bản sạch; ô trống ở bản thô tính là khác.
    \item Trả về \code{brand_id}, \code{so_dong_lech}; sắp \code{brand_id} tăng dần.
  \end{yeucau}
  \goiy{\code{NOT (r.revenue <=> c.revenue)} đúng cả khi \code{r.revenue} là NULL.}
\end{buoc}
\begin{buoc}[ma=b3, diem=8, hinhthuc=write, bang={noise_log, campaigns_raw, campaigns}]{Đối chiếu với nhật ký nhiễu}
  \begin{yeucau}
    \item Nối mỗi dòng nhật ký với đúng dòng của bản thô mà dòng nhật ký đó ghi (cùng \code{brand_id} và \code{row_index}), rồi với chiến dịch ở bản sạch.
    \item Theo \code{noise_type}: số dòng nhật ký, và số dòng mà dòng tương ứng ở bản thô có ít nhất một trong sáu cột số khác bản sạch (cùng quy tắc bước 2).
    \item Trả về \code{noise_type}, \code{so_dong_nhat_ky}, \code{so_dong_lech_chi_so}; sắp \code{noise_type} tăng dần.
  \end{yeucau}
  \goiy{Nối nhật ký với bản thô bằng cả hai cột \code{brand_id} và \code{row_index}; nối theo \code{campaign_id} sẽ khớp cả bản chép của dòng trùng.}
\end{buoc}
\begin{buoc}[ma=b4, diem=7, hinhthuc=ketluan, bang={campaigns_raw, noise_log}]{Kết luận}
  \begin{luachon}
    \phuongan*{Chênh lệch chỉ số đến từ ba loại nhiễu ô trống, số âm, phễu vô lý; dòng trùng chỉ làm thừa dòng, nên dùng được bản sạch.}
      {Đúng: Bước 1 cho thấy phần thừa của bản thô chỉ là dòng trùng; bước 2 và 3 cho thấy số dòng lệch khớp đúng tổng ba loại nhiễu kia.}
    \phuongan{Nhật ký có nhiều dòng hơn số dòng bản thô lệch chỉ số ở bước 2, nên nhật ký không khớp và cần làm lại bản sạch.}
      {Ngày không hợp lệ, khoảng trắng thừa và dòng trùng không chạm sáu cột số, nên không làm lệch chỉ số; bước 3 cho thấy điều này.}
    \phuongan{File của một thương hiệu lệch nhiều hơn hẳn hai file còn lại ở bước 2, nên chỉ cần làm sạch lại file đó.}
      {Bước 1 và 2 cho thấy các file có số dòng thừa và số dòng lệch gần như nhau; nhiễu được chèn đều theo file.}
  \end{luachon}
\end{buoc}
\ghichuthietke{Case Trung bình: 3 bảng chỉ ký hiệu; đường nối \code{row_index} vẽ để gợi khoá dòng nhưng nối đúng cần thêm \code{brand_id} (nêu ở lưu ý và gợi ý bước 3). Trên public: mỗi thương hiệu 50 dòng trùng, 350 dòng lệch chỉ số; bước 3 chỉ \code{'NULL'}, \code{'NEGATIVE'}, \code{'FUNNEL_ERROR'} có dòng lệch. Bài sai: bước 1 đếm mã không DISTINCT; bước 2 dùng \code{<>} bỏ sót ô trống; bước 3 nối theo \code{campaign_id} (dòng trùng nhân đôi) hoặc chỉ theo \code{row_index} (nhân ba). hidden-1 bỏ dòng dữ liệu thô và nhật ký của chiến dịch đuôi 7, nên ba thương hiệu không còn lệch bằng nhau (321/323/315); phương án đúng không dựa vào ý ``lệch như nhau'', phương án sai 3 vẫn sai vì chênh nhau rất ít. Sửa sau kiểm C: rút phương án đúng, cân độ dài ba phương án (K32).}
\end{cauhoi}

\end{document}
